

% Flow-chart styles (panel, ctpnode, description*) are defined in the lambda.tex preamble.
%
% The legend is split into four separate figures so that each panel can be
% cross-referenced individually. Each carries an empty \caption{}, so amsart
% prints just the small "Figure N." label with no descriptive text. The legend
% is \input between two \clearpage commands in lambda.tex, so it has a page to
% itself; the [H] placement keeps the panels in order and the \vfill between
% them spreads the four panels out to fill that page.

\begingroup
% Scale 1 so the legend's nodes match the size of the ctpnodes in the
% flow-chart examples (Figures 5-7), which are drawn unscaled.
\newcommand{\legendscale}{1}%
\setlength{\intextsep}{0.1cm}%
\setlength{\textfloatsep}{0.1cm}%
\setlength{\abovecaptionskip}{3pt}%
\setlength{\belowcaptionskip}{0pt}%

%%%%%%%%%%%%%%%%%%%%%%%%%%%%%%%%%%%%%%%%%%%%%%%%%%%%%%%%%%%%%%
%% PANEL 1 — TERMINAL NODES
%%%%%%%%%%%%%%%%%%%%%%%%%%%%%%%%%%%%%%%%%%%%%%%%%%%%%%%%%%%%%%
\begin{figure}[H]
\centering
\scalebox{\legendscale}{%
\begin{tikzpicture}[>=latex, line join=bevel]
\node (leafpanel) [panel] at (0,0) {
  \begin{tikzpicture}
    % Fix the panel content width so all four panels are the same width
    \node[inner sep=0] at (-7.75,0) {}; \node[inner sep=0] at (7.75,0) {};
    % Title
    \node[font=\bfseries] at (0,2) {Terminal Nodes};

    % Row of nodes
    \node (gray) at (-5,1) [ctpnode, fill=CtpOverlay2] {$\mathbf{(n,s,f)}$};
    \node (peach1) at (0,1) [ctpnode, fill=CtpPeach] {$\mathbf{X(s,f)}$};

    \node (red) at (5,1) [ctpnode, fill=CtpRed] {$\mathbf{(n,s,f)}$};
    % Description boxes
    \node[description] at (-5,-0.2) {$\mathcal{D}_r^{n,s,f} = \textcolor{CtpPeach}{\mathbf{0}} + \textcolor{CtpGreen}{\operatorname{Hom}(\mathcal{U}_r^{n,s,f}, \mathcal{U}_r^{n,s-1,f+r})}$};
    \node[description-small] at (5,-0.2) {Trivial target};
    \node[description-small] at (0,-0.2) {Stable input};
  \end{tikzpicture}
};
\end{tikzpicture}}
\caption{}\label{fig:legend-terminal}
\end{figure}

\vfill

%%%%%%%%%%%%%%%%%%%%%%%%%%%%%%%%%%%%%%%%%%%%%%%%%%%%%%%%%%%%%%
%% PANEL 2 — INTERNAL NODES
%%%%%%%%%%%%%%%%%%%%%%%%%%%%%%%%%%%%%%%%%%%%%%%%%%%%%%%%%%%%%%
\begin{figure}[H]
\centering
\scalebox{\legendscale}{%
\begin{tikzpicture}[>=latex, line join=bevel]
\node (innerpanel) [panel] at (0,0) {
  \begin{tikzpicture}
    % Fix the panel content width so all four panels are the same width
    \node[inner sep=0] at (-7.75,0) {}; \node[inner sep=0] at (7.75,0) {};
    % Title
    \node[font=\bfseries] at (0,2) {Internal Nodes};

    % Row of nodes
    \node (blue) at (-5,1) [ctpnode, fill=CtpSky] {$\mathbf{(n,s,f)}$};
    \node (green) at (0,1) [ctpnode, fill=CtpGreen] {$\mathbf{(n,s,f)}$};
    \node (red2) at (5,1) [ctpnode, fill=CtpRed] {$\mathbf{(n,s,f)}$};

    % Description boxes
    \node[description-wide] at (-5,-0.2) {$\{0\} < \textcolor{CtpGreen}{\mathcal{I}_r^{n, s, f}} < \operatorname{Hom}\big(\mathcal{U}_r^{n, s, f}, \mathcal{U}_r^{n, s - 1, f + r}\big)$};
    \node[description-small] at (0,-0.2) {$\mathcal{D}_r^{n, s, f} = \textcolor{CtpPeach}{\mathbf{d}_r} + \textcolor{CtpGreen}{\{0\}}$};
    \node[description-small] at (5,-0.2) {$\mathcal{D}_r^{n, s, f} = \textcolor{CtpPeach}{\mathbf{0}} + \textcolor{CtpGreen}{\{0\}}$};
  \end{tikzpicture}
};
\end{tikzpicture}}
\caption{}\label{fig:legend-internal}
\end{figure}

\vfill

%%%%%%%%%%%%%%%%%%%%%%%%%%%%%%%%%%%%%%%%%%%%%%%%%%%%%%%%%%%%%%
%% PANEL 3 — ARROWS
%%%%%%%%%%%%%%%%%%%%%%%%%%%%%%%%%%%%%%%%%%%%%%%%%%%%%%%%%%%%%%
\begin{figure}[H]
\centering
\scalebox{\legendscale}{%
\begin{tikzpicture}[>=latex, line join=bevel]
\node (arrowpanel) [panel] at (0,0) {
  \begin{tikzpicture}
    % Fix the panel content width so all four panels are the same width
    \node[inner sep=0] at (-7.75,0) {}; \node[inner sep=0] at (7.75,0) {};
    % Title
    \node[font=\bfseries] at (0,2.8) {Maps};

    % First arrow diagram
    \begin{scope}[shift={(-4,0)}]
      \node (top1) at (0,2) [ctpnode, fill=CtpBase, text=CtpText] {$\mathbf{(n,s,f)}$};
      \node (left1) at (-1.5,1) [ctpnode, fill=CtpBase, text=CtpText] {$\mathbf{F(n,s,f)}$};
      \node (right1) at (1.5,1) [ctpnode, fill=CtpBase, text=CtpText] {$\mathbf{(n,s,f)}$};

      \draw [CtpText,->, very thick] (top1) to[out=225,in=90]
        node [above left,CtpText,inner sep=2pt] {$\mathbf{F}$} (left1);
      \draw [CtpText, very thick] (top1) to[out=315,in=90] (right1);
    \end{scope}

    % Second arrow diagram (copy)
    \begin{scope}[shift={(4,0)}]
      \node (top2) at (0,2) [ctpnode, fill=CtpBase, text=CtpText] {$\mathbf{F(n,s,f)}$};
      \node (left2) at (-1.5,1) [ctpnode, fill=CtpBase, text=CtpText] {$\mathbf{(n, s, f)}$};
      \node (right2) at (1.5,1) [ctpnode, fill=CtpBase, text=CtpText] {$\mathbf{F(n,s,f)}$};

      \draw [CtpText,->, very thick] (left2) to[out=90,in=225]
        node [above left,CtpText,inner sep=2pt] {$\mathbf{F}$} (top2);
      \draw [CtpText, very thick] (top2) to[out=315,in=90] (right2);
    \end{scope}

    % Description boxes for arrow diagrams
    \node[description-leibniz] at (-4,-0.5) {$\mathcal{D}_r^{n, s, f} \update \mathrm{F}_*^{-1}\!\big(\mathrm{F}^*(\mathcal{D}_r^{\mathrm{F}(n, s, f)})\big)$};
    \node[description-leibniz] at (4,-0.5) {$\mathcal{D}_r^{\mathrm{F}(n,s,f)} \update (\mathrm{F}^*)^{-1}\!\big(\mathrm{F}_*(\mathcal{D}_r^{n, s, f})\big)$};
  \end{tikzpicture}
};
\end{tikzpicture}}
\caption{}\label{fig:legend-maps}
\end{figure}

\vfill

%%%%%%%%%%%%%%%%%%%%%%%%%%%%%%%%%%%%%%%%%%%%%%%%%%%%%%%%%%%%%%
%% PANEL 4 — LEIBNIZ RULE
%%%%%%%%%%%%%%%%%%%%%%%%%%%%%%%%%%%%%%%%%%%%%%%%%%%%%%%%%%%%%%
\begin{figure}[H]
\centering
\scalebox{\legendscale}{%
\begin{tikzpicture}[>=latex, line join=bevel]
\node (leibnizpanel) [panel] at (0,0) {
  \begin{tikzpicture}
    % Fix the panel content width so all four panels are the same width
    \node[inner sep=0] at (-7.75,0) {}; \node[inner sep=0] at (7.75,0) {};
    % Title
    \node[font=\bfseries] at (0,3.0) {Products};

    % Left — deducing the left factor ($\mathrm{E}y$ a cycle)
    \begin{scope}[shift={(-4.0,0)}]
      \node (ltop1) at (0,2.2) [ctpnode, fill=CtpBase, text=CtpText] {$\mathbf{(n,s,f)}$};
      \node (lleft1) at (-2.2,1.2) [ctpnode, fill=CtpBase, text=CtpText] {$\mathbf{(n,s,f)}$};
      \node (lmid1) at (0,1.2) [ctpnode, fill=CtpBase, text=CtpText] {$\mathbf{(n{+}s,s',f')}$};
      \node (lright1) at (2.3,1.2) [ctpnode, fill=CtpBase, text=CtpText] {$\mathbf{(n,s{+}s',f{+}f')}$};

      \draw [CtpText, very thick] (ltop1) to[out=270,in=90] (lleft1);
      \draw [CtpText, very thick] (ltop1) to[out=270,in=90] (lmid1);
      \draw [CtpText, very thick] (ltop1) to[out=270,in=90] (lright1);

      \node[description-leibniz] at (0,0.1) {$\mathcal{D}_r^{n,s,f} \update (R_y)_*^{-1}\!\big(R_{\mathrm{E}y}^*(\mathcal{D}_r^{n, s + s', f + f'})\big)$};
    \end{scope}

    % Right — deducing the right factor ($x$ a cycle)
    \begin{scope}[shift={(4.0,0)}]
      \node (ltop2) at (0,2.2) [ctpnode, fill=CtpBase, text=CtpText] {$\mathbf{(n{+}s,s',f')}$};
      \node (lleft2) at (-2.2,1.2) [ctpnode, fill=CtpBase, text=CtpText] {$\mathbf{(n,s,f)}$};
      \node (lmid2) at (0,1.2) [ctpnode, fill=CtpBase, text=CtpText] {$\mathbf{(n{+}s,s',f')}$};
      \node (lright2) at (2.3,1.2) [ctpnode, fill=CtpBase, text=CtpText] {$\mathbf{(n,s{+}s',f{+}f')}$};

      \draw [CtpText, very thick] (ltop2) to[out=270,in=90] (lleft2);
      \draw [CtpText, very thick] (ltop2) to[out=270,in=90] (lmid2);
      \draw [CtpText, very thick] (ltop2) to[out=270,in=90] (lright2);

      \node[description-leibniz] at (0,0.1) {$\mathcal{D}_r^{n+s,s',f'} \update \bigl((L_x)_* \circ \mathrm{E}^*\bigr)^{-1}\!\big((L_x \circ \mathrm{E})^*(\mathcal{D}_r^{n, s + s', f + f'})\big)$};
    \end{scope}

    % Bottom center — deducing the product (full Leibniz rule)
    \begin{scope}[shift={(0,-3.3)}]
      \node (ltop3) at (0,2.2) [ctpnode, fill=CtpBase, text=CtpText] {$\mathbf{(n,s{+}s',f{+}f')}$};
      \node (lleft3) at (-2.2,1.2) [ctpnode, fill=CtpBase, text=CtpText] {$\mathbf{(n,s,f)}$};
      \node (lmid3) at (0,1.2) [ctpnode, fill=CtpBase, text=CtpText] {$\mathbf{(n{+}s,s',f')}$};
      \node (lright3) at (2.3,1.2) [ctpnode, fill=CtpBase, text=CtpText] {$\mathbf{(n,s{+}s',f{+}f')}$};

      \draw [CtpText, very thick] (ltop3) to[out=270,in=90] (lleft3);
      \draw [CtpText, very thick] (ltop3) to[out=270,in=90] (lmid3);
      \draw [CtpText, very thick] (ltop3) to[out=270,in=90] (lright3);

      \node[description, text width=12.2cm] at (0,0.1) {$\mathcal{D}_r^{n, s+s', f+f'} \update \bigl((M \circ (\mathrm{id} \otimes \mathrm{E}))^*\bigr)^{-1}\big(M_*(\mathcal{D}_r^{n, s, f} \otimes \mathrm{id}) + M_*(\mathrm{id} \otimes \mathrm{E}^*(\mathcal{D}_r^{n + s, s', f'}))\big)$};
    \end{scope}
  \end{tikzpicture}
};
\end{tikzpicture}}
\caption{}\label{fig:legend-products}
\end{figure}

\endgroup
